% !PRINTED 6
\weberpage{6}
%<<CONT
und nennen dies eine \gesperrt{lineare Substitution}. Deuten wir $x$ und $x'$ wie im § 1 als Punkte auf zwei geraden Linien $L$ und $L'$, so ist durch (2) eine \gesperrt{gegenseitig eindeutige Beziehung} der Punkte von $L$ und $L'$, eine \gesperrt{Abbildung}, festgelegt. Den Werten $x=\infty$ und $x'=\infty$ entsprechen die unendlich fernen Teile der Geraden, die wir ebenfalls als bestimmte Punkte betrachten. Es entspricht dann der Punkt $x'=\infty$ dem Punkt $x=\alpha:\gamma$ und der Punkt $x=\infty$ dem Punkt $x'=-\delta:\gamma$, und nur wenn $\gamma=0$, die Substitution (2) also \gesperrt{ganz} ist, entsprechen sich die unendlich fernen Punkte auf $L$ und $L'$ gegenseitig.

Die Substitution (2) ändert sich nicht, wenn die vier Transformationszahlen mit demselben Faktor multipliziert werden. Die Determinante $r$ vervielfältigt sich mit dem Quadrate dieses Faktors, und man kann daher diesen Faktor so bestimmen, daß die Determinante einen beliebig gegebenen Wert, z. B. den Wert 1 erhält.

Wendet man die Substitution (2) auf irgend zwei Punkte $x_1, x_2$ und die zugehörigen $x_1', x_2'$ an, so ergibt sich
\[
(x_1x_2) = \frac{r\,(x_1'x_2')}{(\gamma x_1'+\delta)(\gamma x_2'+\delta)},
\]
und wenn man ein zweites Punktepaar $x_3, x_4$ und das entsprechende $x_3', x_4'$ hinzunimmt:
\[
(x_1x_2)(x_3x_4) = \frac{r^2(x_1'x_2')(x_3'x_4')}{(\gamma x_1'+\delta)(\gamma x_2'+\delta)(\gamma x_3'+\delta)(\gamma x_4'+\delta)}.
\]
Vertauscht man hierin $x_2$ mit $x_3$ und bildet den Quotienten, so folgt
\weq{3}{\frac{(x_1x_2)(x_3x_4)}{(x_1x_3)(x_2x_4)} = \frac{(x_1'x_2')(x_3'x_4')}{(x_1'x_3')(x_2'x_4')}.}

Der Ausdruck auf der linken Seite wird das Doppelverhältnis des Punktepaares $x_1x_4$ zu dem Punktepaar $x_2x_3$ genannt, und es ist also in dieser Formel der Satz enthalten:

\gesperrt{Das Doppelverhältnis zweier Punktepaare ist bei gleichzeitiger linearer Transformation der vier Punkte invariant.}

Vier Punkte lassen sich auf sechs Arten in zwei Paare zerlegen. Setzen wir aber
\begin{align*}
a &= (x_2x_3)(x_1x_4) \\
b &= (x_3x_1)(x_2x_4), \\
c &= (x_1x_2)(x_3x_4)
\end{align*}
%CONT>>
